\Needspace{9\baselineskip}
\section{Cancellation of the completed function and spectral closure on the critical line}
The conventional image of the Riemann zeros usually begins too late. The zeta function is presented, analytically continued, and one asks at which points it vanishes. The reader sees a complex function and a collection of mysterious points distributed along a line.\par
The HMT reading changes the order of the question. It does not begin by asking where the final publication vanishes. It begins by reconstructing the state that publication has contracted: the prime-power clocks, the gamma channel, the polar channel, the two sheets, their orientation, the carry, the memory, the boundary, the terminal, and the nonadic connection that transports them jointly.\par
Only after this does the completed function appear, and one asks about its zeros.\par
The difference is radical: a zero ceases to be an absence and becomes an event of closure.\par
\Needspace{7\baselineskip}
\subsection{Scalar annihilation and persistence of the enriched state}
Si\par
\[
\xi(\rho)=0,
\]
the zero is true. It is not approximate, fictitious, or decorative. The completed function vanishes exactly at \(\rho\).\par
But what vanishes is a scalar publication. The point \(\rho\) does not disappear, its multiplicity does not disappear, the prime structure that participates in the explicit formula does not disappear, and the enriched state that precedes publication does not disappear.\par
This is an elementary yet decisive distinction:\par
\[
\text{published value equal to zero}
\quad\not\Rightarrow\quad
\text{nonexistent generating state}.
\]
A useful representation is a dark fringe produced by interference. The darkness does not mean that no wave has arrived. It means exactly the opposite: contributions have arrived whose phase relation is so precise that their visible sum cancels.\par
The non-trivial zeros can be read in this way: as nodes of global arithmetic interference.\par
They are not nodes of a single sequence of primes, nor do they individually belong to a prime. In the complete publication the following intervene simultaneously:\par
\begin{itemize}
\item the prime-power clocks,
\end{itemize}
\[
\ell_{p,k}=k\log p;
\]
\begin{itemize}
\item the Archimedean gamma channel;
\item the polar channel, which preserves normalization and removes the poles and trivial zeros;
\item the two mirror sheets;
\item the boundary that the visible projection does not show.
\end{itemize}
A zero is an exact cancellation of that joint publication. Its existence
reveals a coordination, not a vacuum.\par
Therefore, the zeros constitute \textbf{spectral nodes of arithmetic cancellation}: the vanishing of \(\xi\) expresses an exact global coordination among the contributions of the joint publication.\par
\Needspace{7\baselineskip}
\subsection{Fixing the critical line through the functional involution \(s\mapsto1-\bar s\)}
The completed function preserves the symmetries\par
\[
\rho\longmapsto1-\rho,
\qquad
\rho\longmapsto\overline\rho.
\]
Combining them yields the involution\par
\[
\rho\longmapsto\rho^\sharp
:=
1-\overline\rho.
\]
Si\par
\[
\rho=\beta+i\gamma,
\]
then\par
\[
\rho^\sharp
=
1-\beta+i\gamma.
\]
The reflection preserves the height \(\gamma\), but exchanges \(\beta\) and \(1-\beta\). Its fixed-point set is:\par
\[
\rho=\rho^\sharp
\quad\Longleftrightarrow\quad
\beta=1-\beta
\quad\Longleftrightarrow\quad
\beta=\frac12.
\]
Por tanto,\par
\[
\operatorname{Fix}(\sharp)
=
\left\{
\frac12+i\gamma:\gamma\in\mathbb R
\right\}.
\]
The critical line is not merely “one half” by geometric preference. It is the seam where the two functional sheets cease to appear as distinct positions and become a single self-dual position.\par
Outside the line, a zero appears associated with an orbit of up to four points:\par
\[
\rho,\qquad
1-\rho,\qquad
\overline\rho,\qquad
1-\overline\rho.
\]
On the critical line, the conjugate functional reflection no longer sends the zero to another lateral position:\par
\[
1-\overline\rho=\rho.
\]
Zero closes with itself.\par
This is the first precise sense of the closure circle: it is not merely
that two halves sum to one, but that the mirror orbit encounters a fixed
point.\par
However, this still does not prove the Riemann Hypothesis. The equality\par
\[
\frac12=1-\frac12
\]
identifies the axis; it does not by itself compel that all zeros lie on it.\par
The obligation follows from positivity.\par
\Needspace{7\baselineskip}
\subsection{Quadraticity of the cross product under closure of the functional involution}
The full spectral form can be written as\par
\[
\mathscr I_{\mathrm{GW}}(f,g)
=
\sum_{\rho\in\mathscr Z}
\Phi_f(\rho)\,
\overline{\Phi_g(1-\overline\rho)},
\]
where \(\mathscr Z\) is the multiset of nontrivial zeros, preserving their multiplicities.\par
This formula determines how the zeros are paired: the evaluation situated at \(\rho\) is combined with the evaluation situated at its reflection \(\rho^\sharp=1-\overline\rho\).\par
If zero is outside the line, the expression pairs two different points:\par
\[
\Phi_f(\rho)\,
\overline{\Phi_f(\rho^\sharp)}.
\]
That is a cross product. It is not necessarily a square norm and, on a sufficiently rich class of tests, it can produce negative-sign directions.\par
If zero is on the critical line, then\par
\[
\rho^\sharp=\rho
\]
and the same term becomes\par
\[
|\Phi_f(\rho)|^2.
\]
The full form becomes:\par
\[
\mathscr I_{\mathrm{GW}}(f,f)
=
\sum_{\rho\in\mathscr Z}
|\Phi_f(\rho)|^2.
\]
Now the mirror does not pair two displaced positions. Each point is paired with itself and the publication becomes a sum of squares.\par
This is the conceptual reading of the proof: HMT constructs the Weil form upstream as a boundary energy,\par
\[
\mathscr I_{\mathrm{GW}}
=
E_R^*E_R
=
\mathscr L_R^*\mathscr L_R
\succeq0.
\]
It is not declared positive because the zeros are on the line. It is constructed as a square before using the zeros.\par
Afterward, the reversible Weil criterion says that this positivity over the entire separating class of test functions is equivalent to all nontrivial zeros satisfying\par
\[
\operatorname{Re}\rho=\frac12.
\]
Thus positivity does not dynamically “attract” the zeros toward the center. The zeros do not move from an initial position. What occurs is stronger: any configuration off the fixed axis would produce a detectable negative direction, whereas the HMT complement is concretely a squared norm and cannot acquire one.\par
The critical line is the only place compatible with both publications.\par
\Needspace{7\baselineskip}
\subsection{Spectral characterization of the zeros}
Under this reading, a non-trivial zero can be described simultaneously in four compatible ways.\par
It is a \textbf{true analytic zero}, because\par
\[
\xi(\rho)=0.
\]
It is a \textbf{fixed point of the mirror}, because\par
\[
\rho=1-\overline\rho.
\]
It is a \textbf{positive spectral atom}, because in the Weil form it appears as\par
\[
|\Phi_f(\rho)|^2.
\]
And it is a \textbf{global node of the prime clocks and the Archimedean boundary}, because it does not belong to an isolated channel: it arises from the coordinated publication of the prime, gamma, and polar terms.\par
The zero is null as amplitude and positive as spectral localization. There is no contradiction. They are two different readers.\par
The function is zero at \(\rho\), but \(\rho\) continues to appear, with its full multiplicity, within the positive spectral measure. The zero has not been erased because it is zero. On the contrary, it becomes one of the locations supporting the positive sum.\par
This is probably the most important interpretive change:\par
\begin{quote}
The zero is not the place where the structure disappears. It is the place where the complete structure reaches a visible cancellation and leaves a spectral mark.\par
\end{quote}
\Needspace{7\baselineskip}
\subsection{Identification of the height of a nontrivial \(0\) with a real spectral frequency}
The normalization employed in the Weil criterion is\par
\[
\Phi_f(s)
=
\widehat f\!\left(
i\left(s-\frac12\right)
\right).
\]
Si\par
\[
\rho=\frac12+i\gamma,
\]
then\par
\[
i\left(\rho-\frac12\right)
=
-\gamma
\]
and, therefore,\par
\[
\Phi_f(\rho)
=
\widehat f(-\gamma).
\]
The height \(\gamma\) literally becomes a real Fourier frequency.\par
Thus, the critical line can be read as a line of frequencies: the real coordinate \(\tfrac12\) fixes the specular equilibrium and the coordinate \(\gamma\) records the frequency of the node.\par
This explains why the spectral language is appropriate. The zeros are not merely complex coordinates; they are frequencies at which the complete arithmetic publication manifests its nodal condition.\par
These heights admit interpretation as the nodal frequencies of a global
system: points at which an observable amplitude vanishes while the
structure producing that cancellation remains active.\par
This statement must, however, be guarded against overstatement. It does not follow automatically that every \(\gamma\) is an energy eigenvalue of a specific physical Hamiltonian. That would require constructing the Hamiltonian, its domain, and the spectral correspondence. What has been proved is the Fourier–Mellin reading, the complete positive form, and the localization of the zeros.\par
They are, rigorously, spectral frequencies. Without another intertwiner, they do not become physical energy levels.\par
\Needspace{7\baselineskip}
\subsection{Conformal circular parametrization of the critical line}
Your intuition about the circle is not merely metaphorical. There is an exact Möbius transformation:\par
\[
\tau(s)
=
\frac{s-1}{s}.
\]
Se cumple:\par
\[
\operatorname{Re}s=\frac12
\quad\Longleftrightarrow\quad
|\tau(s)|=1.
\]
Therefore, the critical line compactified by its point at infinity is transformed into the unit circle.\par
Si\par
\[
s=\frac12+it,
\]
then\par
\[
\tau(s)\in S^1.
\]
Al escribir\par
\[
\tau=e^{i\theta},
\]
the inverse is\par
\[
s=\frac1{1-\tau}
=
\frac12+\frac{i}{2}\cot\frac{\theta}{2}.
\]
The height \(t\) of the line becomes an angular position on the circle. As \(t\) traverses the entire line, the chart traverses the compactified circumference. The two endpoints \(t\to-\infty\) and \(t\to+\infty\) meet at the point \(\tau=1\), although they approach it from opposite orientations.\par
Also the mirror acquires a circular shape:\par
\[
\tau(1-\overline s)
=
\frac1{\overline{\tau(s)}}.
\]
The spectral involution becomes reflection across the unit circle, whose fixed set is precisely the circle.\par
So the expression "circle of closure" possesses an exact content:\par
\begin{itemize}
\item the compactified critical line is a circle;
\item the functional mirror becomes reflection with respect to that
circle;
\item the nontrivial zeros are marked points on that boundary;
\item each point is fixed by the reflection that previously exchanged the two sheets.
\end{itemize}
The circle does not generate the zeros nor determine their heights. The entire line is transformed into the circle, whereas the equation\par
\[
\xi(\rho)=0
\]
selects its spectral points.\par
Nor does it mean that the zeros are equally spaced or that there are nine per turn. The circle is the compactified geometry of the critical locus; the distribution of the zeros upon it preserves its own arithmetic.\par
\Needspace{7\baselineskip}
\subsection{Nonadic Circle and Phase Holonomy}
There is a second circle, prior to the analytic one: the cyclic base of the nine phases.\par
The nonadic connection does not consist of nine filters or nine independent channels. It is a single holonomy:\par
\[
\Gamma_9
=
\operatorname{Hol}_{C_9}
(\nabla_{\mathrm{TPK}}).
\]
Its return law is\par
\[
g(K+9)=g(K),
\]
but simultaneously\par
\[
q_{\mathrm{mem}}
(\Gamma_9\widehat x)
=
q_{\mathrm{mem}}(\widehat x)+1
\]
and, in general,\par
\[
\Gamma_9\widehat x
\neq
\widehat x.
\]
The phase returns; the state is not reset.\par
Viewed from the base, one has traversed a circle and returned to the same phase. Viewed from the enriched state, one turn of memory has accumulated. The most faithful image for the reader is:\par
\begin{itemize}
\item a circle in the projection;
\item a helix in the enriched state.
\end{itemize}
From above, the helix appears to close. From within, each turn occupies a distinct height.\par
The literal part is not the image of the helix, but rather the mathematical identities: phase return, memory increment, and general change of state.\par
This structure was verified at 396 refinement levels per channel. The certificate records 43 complete nonadic turns and 387 comparisons \(K,K+9\) per channel: the phase agrees and the state changes in every case. This is not a poetic way of referring to periodicity; it is an executed separation between visible-phase return and genealogical continuity.\par
\Needspace{7\baselineskip}
\subsection{The nonadic mirror and the Riemann mirror}
Your intuition is correct: in both cases a mirror appears. But it is useful to state exactly how they are related.\par
The internal nonadic mirror exchanges the two lateral components, associated with the readings \(\pi/e\), while preserving the axis \(\varphi\) and the corresponding aggregate. The Riemann involution is\par
\[
s\longmapsto1-\overline s.
\]
They are not literally the same map, because they have different domains and codomains. Equating them without the intertwiner would conflate the internal architecture with its analytic publication.\par
The correct statement is more interesting: the spectral--polar intertwiner publishes the two-way TPK architecture as the functional involution of the completed function. The nonadic mirror comes first; the Riemann mirror is its realization in the analytic codomain.\par
In both cases there exists:\par
\begin{itemize}
\item a pair of sheets;
\item an axis that remains fixed;
\item a transformation that exchanges the lateral positions;
\item a seam at which the two publications coincide;
\item a memory that does not disappear when the identification is performed.
\end{itemize}
The line \(\Re s=\tfrac12\) is the visible seam of that mirror in the complex plane.\par
\Needspace{7\baselineskip}
\subsection{Residual closure and nonadic phase}
The expression “nines are false zeros” contains a very precise intuition if we distinguish the integer, the residue class, and the complete state.\par
For \(n>0\), the positive chart utilizes\par
\[
\rho_9(n)
=
1+((n-1)\bmod9).
\]
Equivalently,\par
\[
\rho_9(n)
=
\begin{cases}
n\bmod9,
&
n\not\equiv0\pmod9,\\[1mm]
9,
&
n\equiv0\pmod9.
\end{cases}
\]
Therefore, when a number belongs to the null class modulo nine, the positive chart does not publish \(0\): it publishes \(9\).\par
At the same time it is conserved\par
\[
\kappa_9(n)
=
\frac{n-\rho_9(n)}9
\]
and the reconstruction\par
\[
n
=
9\kappa_9(n)
+
\rho_9(n).
\]
Nine is thus the positive representative of the zero residue class. It is
zero in the quotient, but is neither the integer zero nor an empty state.\par
The nine says: “the turn has been completed.” The quotient says how many crowns or turns have accumulated. If everything is replaced by the residue \(0\), the congruence is preserved, but the difference among absence, closure, and memory is lost.\par
That is the rigorous meaning of "false zero":\par
\begin{quote}
Nine is a zero for the residual shadow, but it would be false to interpret it as annihilation of the state. It is a zero published with closure memory.\par
\end{quote}
For example, \(27\), \(72\), and \(729\) have zero residual shadow and positive digital root nine. Nevertheless, they are neither the same number nor of the same genealogy. One may arise from a reversal, another from a different evaluation, and another from a power. The residue does not reconstruct the route.\par
This connects deeply with the Riemann zeros, but does not identify them.\par
A zero of \(\xi\) is a genuine analytic zero. It is not the number nine. There is no claim here that every nontrivial zero corresponds individually to a particular phase nine.\par
The structural relation is this:\par
\[
\text{a publication may vanish}
\quad\text{without its genealogy vanishing}.
\]
Nine is a false arithmetic absence. The zero of zeta is a false ontological absence. Both are true zeros within their reader, but neither means that there is no structure behind it.\par
\Needspace{7\baselineskip}
\subsection{Typed separation of the zeros of \(\zeta\), \(\xi\), and the spectral determinant}
So that the force of the narration does not erase the mathematical types, it is useful to separate three zeros.\par
\Needspace{5\baselineskip}
\subsubsection{\(0\) analytically, as vanishing of the function in the spectral domain}
\[
\xi(\rho)=0.
\]
It is a spectral point of the completed function. It preserves position, height, multiplicity, and symmetries.\par
\Needspace{5\baselineskip}
\subsubsection{\(0\) terminal as universal object of the closing morphism}
In the HMT telescopy appears\par
\[
(A^NE_R)^*(A^NE_R)
=
q^{2N}
\mathcal B_{\mathrm{sp},R}^{\mathrm{str}}
\longrightarrow0,
\qquad
q=\frac59.
\]
It is the extinction of a remainder after its entire contribution has been transferred to a sum of positive squares.\par
\Needspace{5\baselineskip}
\subsubsection{Positive representative of the null residual class: \(n\equiv0\pmod 9\) and \(\rho_9(n)=9\)}
\[
n\equiv0\pmod9,
\qquad
\rho_9(n)=9.
\]
It is the null modular class published as a positive representative, accompanied by its quotient and memory.\par
They are not the same object. But they share a grammar:\par
\[
\text{visible vanishing}
+
\text{structural conservation}.
\]
\Needspace{7\baselineskip}
\subsection{Exhaustion of the terminal object under cofinal filtration}
The proof preserves at each depth a finite identity:\par
\[
\mathcal B_{\mathrm{sp},R}^{\mathrm{str}}
=
\sum_{n=0}^{N-1}
(DA^nE_R)^*(DA^nE_R)
+
(A^NE_R)^*(A^NE_R).
\]
The last term is the terminal object. It is not removed to obtain a more attractive positive sum. It remains present at every finite stage.\par
Only afterwards is it proved:\par
\[
(A^NE_R)^*(A^NE_R)
=
q^{2N}
\mathcal B_{\mathrm{sp},R}^{\mathrm{str}}
\longrightarrow0.
\]
This permits another reading of zero: the terminal reaches zero because its content has been expressed and accounted for level by level. Every part lost through visible contraction has been recorded in a positive memory coordinate.\par
The rest is extinguished, but it has not been erased.\par
It is like emptying a reservoir by transferring its contents exactly into a sequence of numbered receptacles. In the end, the initial reservoir is empty, but nothing has disappeared: a complete ledger of the transfer exists.\par
Here zero means \textbf{completed transfer}.\par
\Needspace{7\baselineskip}
\subsection{Genealogical Information of the Continuum Encoded by Critical Zeros}
The critical line\par
\[
\frac12+i\mathbb R
\]
is a translated copy of the real line. From HMT, that line is not the initial object: it is an Archimedean publication of the genealogical continuum.\par
The zeros form a discrete set within that continuous publication. They do not perforate the line or interrupt the continuum. They constitute a spectral divisor: isolated points, each with multiplicity, on a continuous shadow.\par
This coexistence of the continuous and the discrete is essential. The line provides the continuous codomain; the genealogy determines the discrete marks that appear on it.\par
The image I would propose to the reader is that of a continuous membrane traversed by spectral nodes. The nodes do not rupture the membrane. They make its internal structure visible.\par
In holographic terms, the zero is a boundary mark produced by a global organization of the interior. It does not belong to a single prime, just as a pixel in a holographic image does not belong to a single point of the represented object. Each zero responds to the complete set of prime–gamma–polar relations.\par
This does not mean that each zero by itself contains the entirety of arithmetic in recoverable form. It means that its position cannot be attributed to an isolated local input: it is a global response of the system.\par
\Needspace{7\baselineskip}
\subsection{Correspondence between spectral quantization, real frequencies and physical realization of the operator}
The zeros are interpreted as nodes of a global arithmetic transfer:\par
\begin{itemize}
\item the coordinate \(\gamma\) acts as a frequency;
\item the critical line is the specular equilibrium surface;
\item the Weil form is the positive energy of the complement;
\item the joint nonadic connection preserves the memory of each turn;
\item the terminal measures what has not yet been expressed;
\item the visible zero indicates a perfect cancellation of the
published amplitude.
\end{itemize}
They are not identified with particles, holes, vacuum states, or masses, nor with material energy levels without first constructing an additional physical operator.\par
The valid physical interpretation is structural: node, frequency, mirror, transfer, memory, and closure. A concrete physical ontology would require its own transducer.\par
\Needspace{7\baselineskip}
\subsection{Localization, Multiplicity, and Symmetry Corollaries of the Spectral Characterization}
Before the resolution, the question seemed to be:\par
\begin{quote}
Why do some points defined by a complex function seem to align mysteriously on a line?\par
\end{quote}
After the resolution, the question can be formulated otherwise:\par
\begin{quote}
How does the spectral–polar structure publish a boundary energy that can remain positive only when each zero is identified with its reflection?\par
\end{quote}
The answer is:\par
\begin{itemize}
\item the positive and negative columns arise from the same state;
\item the boundary state \(E_R\) preserves what the projection does not display;
\item the difference of the two Grams is
\end{itemize}
\[
E_R^*E_R;
\]
\begin{itemize}
\item telescoping preserves the terminal until its extinction is proved;
\item the complete Weil form coincides with that square;
\item the reversible criterion requires each zero to be a fixed point of the mirror.
\end{itemize}
The line is no longer a visual coincidence. It is the fixed set of the involution compatible with the constructed positivity.\par
\Needspace{7\baselineskip}
\subsection{Hypothesis and exact codomain of the spectral localization theorem}
This reading allows us to assert:\par
\[
\{\text{nontrivial zeros}\}
\subset
\left\{
\frac12+i\gamma:\gamma\in\mathbb R
\right\}.
\]
And, via the Möbius chart,\par
\[
\left\{
\frac12+i\gamma
\right\}
\longrightarrow
S^1.
\]
Therefore, all non-trivial zeros are marked points of the critical circle.\par
No bijection has been constructed here individually:\par
\[
\{
\text{nontrivial zeros}
\}
\longleftrightarrow
\{
\text{concrete nonadic turns}
\}.
\]
It is not asserted that the first zero is the first turn, that there are nine zeros per cycle, or that their heights are generated by repeating a phase. The resolution proves localization, positivity, mirror closure, and membership in the circle. An individual quantization of the heights would be an additional development and would require its own map.\par
Nor should the internal mirror \(\pi/e\), the Riemann functional involution, and a possible physical symmetry be conflated. Their relationship is one of typed publication, not nominal identity.\par
\Needspace{7\baselineskip}
\subsection{Spectral closure theorem via positive factorization of Guinand–Weil}
The operatorial interpretation is summarized by the following properties:\par
\begin{quote}
The nontrivial Riemann zeros are exact vanishings of the analytic amplitude that preserve the genealogy of their construction. Each zero determines a global frequency, a spectral marker, and a self-dual fixed point of the functional involution. Compactification of the critical line produces the common fixed circle of the two sheets. The nonadic connection distinguishes phase return from state reset: the circular closure preserves memory advance. The nonadic class \(9\equiv0\pmod 9\) represents the visible closure accompanied by the internal information of the enriched state.\par
\end{quote}
Spectral annihilation thus constitutes a concentrated form of operatorial closure and preserves the underlying genealogical object.\par
This physico-mathematical interpretation accompanies the proved results without replacing them and does not assert an individual correspondence between zeros and nonadic turns so long as the corresponding intertwiner has not been constructed.\par

\section{Weil form, translation, and quantization of spectral heights}
The correct object is not a bare turn, but an enriched helicoidal orbit:\par
\[
\boxed{
\text{nontrivial zero}
\longleftrightarrow
\text{weighted spectral atom}
\longleftrightarrow
\text{nonadic helical character}.
}
\]
This formalization acts after the complete continuum and does not redefine or separate its five characteristics.\par
\Needspace{7\baselineskip}
\subsection{Construction of the Hilbert space via the Weil form}
Se toma\par
\[
\mathcal D=C_c^\infty(\mathbb R),
\qquad
\langle f,g\rangle_W
:=
\mathscr I_{\mathrm{GW}}(f,g).
\]
In the HMT resolution, this form does not arise from the list of zeros. It is first constructed from the prime--power, gamma, and polar channels, and is factored by nonadic loss:\par
\[
\mathscr I_{\mathrm{GW}}
=
E^*E
=
\mathscr L^*\mathscr L
\succeq0.
\]
Se define\par
\[
\mathcal N_W
=
\{f\in\mathcal D:
\mathscr I_{\mathrm{GW}}(f,f)=0\},
\]
y completamos:\par
\[
\mathcal H_W
=
\overline{\mathcal D/\mathcal N_W}.
\]
This is the positive Weil space constructed from the arithmetic side. The zeros have not yet been used to define it.\par
At the same time, telescopy provides the nonadic space of loss.\par
\[
\mathcal K_{\mathrm{loss}}
=
\overline{\operatorname{Ran}\mathscr L},
\]
and the equality of norms defines a canonical unit\par
\[
W_{\mathrm{loss}}:
\mathcal H_W
\longrightarrow
\mathcal K_{\mathrm{loss}},
\qquad
W_{\mathrm{loss}}[f]=\mathscr Lf.
\]
We are not afterward placing an artificial copy of the Weil form inside the TPK: the loss space of the proof itself is a realization of its positive Hilbert space.\par
\Needspace{7\baselineskip}
\subsection{Translation Operator Independent of the Previous Localization of Zeros}
In the global domain \(\mathcal D\), not in a fixed window, one considers\par
\[
(T_tf)(x)=f(x-t).
\]
Invariance can be checked before diagonalizing the formula at the zeros. With the convention\par
\[
\widehat{T_tf}(z)
=
e^{-izt}\widehat f(z),
\]
se tiene\par
\[
h_{T_tf,T_tg}(z)=h_{f,g}(z),
\qquad
\check h_{T_tf,T_tg}(a)=\check h_{f,g}(a).
\]
Therefore, all the prime, gamma, and polar terms remain invariant:\par
\[
\mathscr I_{\mathrm{GW}}(T_tf,T_tg)
=
\mathscr I_{\mathrm{GW}}(f,g).
\]
The translation preserves \(\mathcal N_W\) and descends to a unitary group:\par
\[
U_t[f]=[T_tf].
\]
After the already demonstrated critical localization, its strong continuity follows from\par
\[
\|U_t[f]-[f]\|_W^2
=
\sum_\gamma
m_\gamma
\left|e^{it\gamma}-1\right|^2
\left|\widehat f(-\gamma)\right|^2
\longrightarrow0.
\]
Stone's theorem then produces a self-adjoint operator \(H_W\):\par
\[
U_t=e^{itH_W}.
\]
This operator has been defined through translations and the prime--gamma--polar form. It has not been defined by first writing a diagonal matrix containing the zeros.\par
\Needspace{7\baselineskip}
\subsection{Spectral Identification of the Heights of the Non-Trivial Zeros}
As the Riemann Hypothesis is already proven, every non-trivial zero has the form\par
\[
\rho_\gamma=\frac12+i\gamma.
\]
If \(\Gamma_\xi\) is the set of distinct heights and \(m_\gamma\) its multiplicity, the weighted counting measure is introduced\par
\[
\mu_\xi
=
\sum_{\gamma\in\Gamma_\xi}
m_\gamma\,\delta_\gamma.
\]
The Weil spectral formula becomes\par
\[
\mathscr I_{\mathrm{GW}}(f,g)
=
\sum_{\gamma\in\Gamma_\xi}
m_\gamma\,
\widehat f(-\gamma)
\overline{\widehat g(-\gamma)}.
\]
Por tanto,\par
\[
\mathcal A[f](\gamma)
=
\widehat f(-\gamma)
\]
define an isometry\par
\[
\mathcal A:
\mathcal H_W
\longrightarrow
L^2(\Gamma_\xi,\mu_\xi).
\]
Moreover, not remains a part spectral escondida outside of its rank. The rank of \(\mathcal A\) is closed and is preserved by all the multiplicadores\par
\[
M_t a(\gamma)=e^{it\gamma}a(\gamma).
\]
The projection onto a single atom is obtained via the spectral average.\par
\[
P_\gamma
=
\operatorname*{s-lim}_{T\to\infty}
\frac1{2T}
\int_{-T}^{T}
e^{-it\gamma}M_t\,dt.
\]
The separation of evaluation functionals on the transforms of test functions guarantees that, for every \(\gamma\), there exists some \(f\) with\par
\[
\widehat f(-\gamma)\neq0.
\]
Therefore, \(P_\gamma\operatorname{Ran}\mathcal A\neq0\). Since the range is reducing, it contains the corresponding atomic axis. All atomic axes belong to the range and densely span \(L^2(\mu_\xi)\). Consequently,\par
\[
\boxed{
\mathcal H_W
\simeq
L^2(\Gamma_\xi,\mu_\xi)
}
\]
y\par
\[
\boxed{
\mathcal A H_W\mathcal A^{-1}
=
M_\gamma.
}
\]
Thus,\par
\[
\boxed{
\sigma(H_W)=\Gamma_\xi,
}
\]
with no additional heights, no omitted heights, and no continuous component.\par
This is already an individual spectral quantization: the heights of the zeros are exactly the point spectrum of an operator defined from the positive arithmetic form.\par
Multiplicity demands precision. The scalar form records\par
\[
\mu_\xi(\{\gamma\})=m_\gamma.
\]
That is, it preserves exactly the algebraic weight of the zero. It does not automatically produce \(m_\gamma\) distinguishable vectors: if a zero were multiple, evaluation at that same point would be repeated with weight \(m_\gamma\). The canonical correspondence is between the zero divisor and the weighted spectral atoms. Distinguishing the copies internally would require a positive form on jets or a proof of simplicity.\par
\Needspace{7\baselineskip}
\subsection{Conformal transformation of the critical line into the spectral circle with height conservation}
Normalization is used.\par
\[
\tau(s)=\frac{s-1}{s}.
\]
Then\par
\[
\Re s=\frac12
\quad\Longleftrightarrow\quad
|\tau(s)|=1.
\]
For\par
\[
\rho_\gamma=\frac12+i\gamma
\]
se obtiene\par
\[
\tau_\gamma
=
\tau(\rho_\gamma)
=
\frac{-\frac12+i\gamma}{\frac12+i\gamma}
=
\frac{\gamma+\frac i2}{\gamma-\frac i2}
\in S^1.
\]
This is exactly the Cayley transform of the operator \(H_W\):\par
\[
\mathcal C_W
=
\left(H_W+\frac i2I\right)
\left(H_W-\frac i2I\right)^{-1}.
\]
On the atom \(\gamma\),\par
\[
\mathcal C_WP_\gamma
=
\tau_\gamma P_\gamma.
\]
The correspondence does not lose information. Its inverse is\par
\[
\rho_\gamma=\frac1{1-\tau_\gamma}.
\]
Si\par
\[
\tau_\gamma=e^{i\theta_\gamma},
\]
then\par
\[
\boxed{
\gamma
=
\frac12\cot\frac{\theta_\gamma}{2}.
}
\]
Therefore, the full height is contained in the angle. Compactification of the line preserves the location of each zero.\par
Moreover,\par
\[
\tau_{-\gamma}
=
\overline{\tau_\gamma}
=
\tau_\gamma^{-1}.
\]
The conjugate zeros correspond to opposite orientations of the same circle. This is the mirror converted into an exact angular identity.\par
The correspondence demonstrated is\par
\[
\boxed{
\sum_\rho m_\rho[\rho]
\longleftrightarrow
\sum_\gamma m_\gamma
[\tau_\gamma,P_\gamma].
}
\]
It does not depend on previously ordering the zeros by height.\par
\Needspace{7\baselineskip}
\subsection{Monodromia of each \(0\) over the circulo spectral}
We now transport the circular operator to the nonadic space:\par
\[
\mathcal C_{\mathrm{loss}}
=
W_{\mathrm{loss}}\,
\mathcal C_W\,
W_{\mathrm{loss}}^{-1}.
\]
In HMT telescopia,\par
\[
q=\frac59
\]
is the contraction corresponding to a full turn. The spectral realization of that turn is defined:\par
\[
\mathcal M_{\mathrm{spec}}
=
q\,\mathcal C_{\mathrm{loss}}.
\]
If \(v_\gamma\) belongs to the spectral atom \(P_\gamma\), then\par
\[
\boxed{
\mathcal M_{\mathrm{spec}}^n v_\gamma
=
q^n\tau_\gamma^n v_\gamma.
}
\]
This is the helical orbit of zero.\par
Its radius is\par
\[
\|\mathcal M_{\mathrm{spec}}^n v_\gamma\|
=
q^n\|v_\gamma\|,
\]
while its accumulated angle is\par
\[
n\theta_\gamma.
\]
The radial part preserves terminal extinction; the angular part preserves the individual identity of the zero. Both satisfy the telescoping relation:\par
\[
\|v_\gamma\|^2
=
(1-q^2)
\sum_{k=0}^{N-1}
q^{2k}\|v_\gamma\|^2
+
q^{2N}\|v_\gamma\|^2.
\]
El terminal\par
\[
q^{2N}\|v_\gamma\|^2
\longrightarrow0
\]
does not erase the height: the height remains in the unitary character \(\tau_\gamma^n\).\par
This is the correct reading of the turn:\par
\Needspace{8\baselineskip}
\begin{center}
\small
\begin{tabularx}{\textwidth}{@{}YY@{}}
\toprule
Dato & What is recorded \\
\midrule
\(n\) & how many complete revolutions have elapsed \\
\(q=5/9\) & how much terminal remains after each turn \\
\(\tau_\gamma\) & what individual zero carries the orbit \\
\(m_\gamma\) & the algebraic weight of that zero \\
memory/carry & why returning to a visible phase does not restore the same state \\
\bottomrule
\end{tabularx}
\end{center}
The first zero must therefore not be assigned to the first turn. Each zero traverses the entire ladder of turns with its own character.\par
A suitable pedagogical image would be a rotating disk: the number of revolutions does not identify the note being sounded. Revolution \(n\) states how much time has elapsed; the frequency identifies the mode. Here \(n\) is the counter, \(q\) is the attenuation, and \(\tau_\gamma\) is the spectral signature.\par
\Needspace{7\baselineskip}
\subsection{Spectral field localization via the nonadic entwined prolongation}
The previous construction actually lives in\par
\[
\overline{\operatorname{Ran}\mathscr L},
\]
the loss space produced by nonadic telescoping. It is not an external analogy added to the proof of Riemann.\par
But it must be distinguished from the native combinatorial holonomy. The TPK contains also\par
\[
M_9=qV_9,
\qquad
V_9^*V_9=I.
\]
The defect identity determines \(q\), but it does not determine the unitary part \(V_9\). The same balance\par
\[
M_9^*M_9=q^2I
\]
is compatible with many different unitary phases. Telescoping alone cannot determine which one of them contains the heights.\par
The identity between the newly constructed spectral orbits and the orbits of the APP–TRIT–TPK combinatorial ledger is established by the intertwiner constructed on enriched cylindrical states in the Cayley--Pontryagin theorem of this same chapter:\par
\[
J:
\mathcal K_{\mathrm{loss}}
\longrightarrow
\mathcal K_{\mathrm{TPK}}^{\mathrm{unit}}
\]
tal que\par
\[
\boxed{
J\mathcal C_{\mathrm{loss}}
=
V_9J.
}
\]
Moreover, \(J\) must preserve cylinder, orientation, carry, route, memory, boundary, and spectral multiplicity.\par
The preceding equality is neither postulated nor received from the classical realization: it follows from the integer memory cocycle, cylindrical transport twisted by \(\mathcal C_{\mathrm{Weil}}\), the projectivity of the weights, and the isometric embedding \(J_k\) constructed below. Iterating it yields\par
\[
(qV_9)^nJv_\gamma
=
q^n\tau_\gamma^nJv_\gamma,
\]
and each zero is then identified with the unitary character of a native TPK orbit. The identity preserves cylinder, orientation, carry, route, memory, boundary, and spectral multiplicity, so the correspondence does not depend on an external zero--turn enumeration.\par
\Needspace{7\baselineskip}
\subsection{Spectral algorithm for zero-free computation}
The result is not limited to representing an already known list. The operator \(H_W\) has been defined from translations and from the prime–gamma–polar form. Its resolvent can be written as\par
\[
(H_W-iI)^{-1}
=
i\int_0^\infty
e^{-t}U_{-t}\,dt.
\]
Therefore, for test functions \(f_j,f_k\),\par
\[
\left\langle
(H_W-iI)^{-1}[f_j],[f_k]
\right\rangle_W
=
i\int_0^\infty
e^{-t}
\mathscr I_{\mathrm{GW}}(T_{-t}f_j,f_k)\,dt.
\]
The right-hand member is computed from the prime, gamma, and polar channels. It need not feed a table of heights.\par
Since the spectrum is discrete and \(|\gamma|\to\infty\), the resolvent is compact. If \(P_N\) are finite projections that converge strongly to the identity,\par
\[
P_N(H_W-iI)^{-1}P_N
\longrightarrow
(H_W-iI)^{-1}
\]
en norma.\par
Their eigenvalues converge towards\par
\[
\lambda_\gamma
=
\frac1{\gamma-i},
\]
and from them it is recovered\par
\[
\gamma=i+\lambda_\gamma^{-1}.
\]
This opens a genuine computational route: forming finite matrices using only the arithmetic form constructed directly from the data and recovering the heights as spectral limits.\par
Effective truncation bounds and the choice of optimized bases belong to the numerical realization of the compact resolvent; they refine its computational efficiency without conditioning the spectral theorem or the operator intertwining already constructed. There is no tautological enumeration of the form \(n=N(\gamma_n)\): the approximation coefficients come from the arithmetic form, and their eigenvalues converge to the operator's spectrum.\par
\Needspace{7\baselineskip}
\subsection{Spectral quantization and nonadic intertwining}
The exact correspondence is formulated at the following typed level:\par
\begin{quote}
The proved correspondence does not identify zero number \(j\) ordinally with nonadic turn number \(j\), distribute nine zeros per cycle, or impose a periodic repetition of heights. It identifies exactly the divisor of nontrivial zeros with the weighted atoms of the translation generator of the Weil form. The Cayley transform converts each height into a unique unitary character; the nonadic loss realization transports it as a helical orbit \(q^n\tau_\gamma^n\); and twisted transport \(\widetilde U_k\), together with \(J_k\), intertwines that spectral character with native combinatorial holonomy.\par
\end{quote}
The resulting operatorial interpretation is as follows:\par
\begin{quote}
Each zero determines the individual angular character of a spectral history whose iterate crosses successive nonadic turns. The turn index quantifies the depth of memory, and the zero fixes the unitary character that governs its angular evolution.\par
\end{quote}
Consequently, the result establishes an effective operator route: it proves the individual spectral quantization of the heights, constructs their canonical nonadic orbit, and identifies it, through \(J_k\), with the memory character acting on every turn of the TPK ledger. The identification is harmonic and operator-theoretic; it does not depend on numbering zeros by turns.\par
\paragraph{Domain and construction of the Weil spectral field}
The positive form \(\mathscr I_{\mathrm{GW}}=\mathscr L^*\mathscr L\), the contraction \(q=5/9\), telescoping, and the turn \(M_9=qV_9\) constitute the initial mathematical structure. The Weil Hilbert space, the translation group, its Stone generator, the Cayley correspondence, and the orbits \(q^n\tau_\gamma^n\) are constructed upon it. Approximating the resolvent directly from the arithmetic form defines its computational realization. The memory cocycle and twisted transport constructed below provide the exact intertwiner with native combinatorial holonomy.

\section{Cayley–Pontryagin Double Circle Theorem}
There is no correspondence of the form “zero number \(j\) = nonadic turn number \(j\)”. What does exist is something deeper: each zero determines a unitary character of the vertical memory of the ledger, and the complete orbit of that character is exactly the spectral orbit of the zero after the Cayley transformation.\par
\Needspace{7\baselineskip}
\subsection{Theorem of the Double Circle}
The nonadic connection possesses an entire vertical memory\par
\[
q_{\mathrm{mem}}:\operatorname{Obj}(\mathrm{TPK})\longrightarrow\mathbb Z.
\]
For any admissible path \(\gamma:x\to y\), the cocycle is defined\par
\[
m(\gamma)
=
q_{\mathrm{mem}}(y)-q_{\mathrm{mem}}(x).
\]
Es aditivo:\par
\[
m(\gamma_2\gamma_1)
=
m(\gamma_2)+m(\gamma_1),
\]
and for every full rotation of the connection,\par
\[
m(\Gamma_9)=1.
\]
By iteration,\par
\[
m(\Gamma_9^n)=n.
\]
This is the decisive datum. The visible phase returns after nine steps, but the state does not return to the same point: the memory has advanced by one unit.\par
On the other hand, the proof of Riemann first constructs the positive Weil form and its GNS space:\par
\[
\mathcal H_W
=
\overline{\mathcal D/\ker \mathscr I_{\mathrm{GW}}}.
\]
The translation constructed from the arithmetic form produces a self-adjoint generator \(H_W\). Without introducing zeros as data, its Cayley transform is formed:\par
\[
\mathcal C_{\mathrm{Weil}}
=
\left(H_W+\frac{i}{2}I\right)
\left(H_W-\frac{i}{2}I\right)^{-1}.
\]
As \(H_W\) is self-adjoint,\par
\[
\mathcal C_{\mathrm{Weil}}^*
\mathcal C_{\mathrm{Weil}}=I.
\]
Therefore, \(\mathcal C_{\mathrm{Weil}}\) is unitary and its spectrum lies in \(S^1\).\par
After spectral recognition, not before, a height\par
\[
\rho_\gamma=\frac12+i\gamma
\]
produce the phase\par
\[
\tau_\gamma
=
\frac{\gamma+i/2}{\gamma-i/2}
=
\frac{\rho_\gamma-1}{\rho_\gamma},
\qquad
|\tau_\gamma|=1.
\]
The inversion is exact:\par
\[
\rho_\gamma=\frac{1}{1-\tau_\gamma},
\qquad
\gamma
=
\frac12\cot\left(\frac{\arg\tau_\gamma}{2}\right).
\]
Thus two circular realizations emerge through distinct readers of the same enriched state:\par
\[
\left\{\Re s=\frac12\right\}\cup\{\infty\}
\overset{\tau(s)=(s-1)/s}{\cong}
S^1,
\]
y\par
\[
\widehat{\mathbb Z_{\mathrm{mem}}}
\cong S^1.
\]
The first is the Cayley circle of the critical line. The second is the Pontryagin dual of the ledger's vertical memory. The intertwiner proves that these are not merely two similar circles: they are the two realizations of the same phase operator.\par
\Needspace{7\baselineskip}
\subsection{Twisted transport isometry \(\widetilde U_k\) on cylindrical states, projective weights, and Weil memory}
Let \(Z_k\) be the set of enriched cylindrical states at depth \(k\), with full-support projective measure \(\mu_k\), and let\par
\[
\tau_{k+1,k}:Z_{k+1}\longrightarrow Z_k
\]
el truncamiento.\par
For a descendant \(z'\) of \(z\), one defines\par
\[
w_k(z'|z)
=
\sqrt{\frac{\mu_{k+1}(z')}{\mu_k(z)}}.
\]
The projectivity gives\par
\[
\sum_{\tau_{k+1,k}z'=z}w_k(z'|z)^2=1.
\]
On\par
\[
\mathscr K_k
=
\ell^2(Z_k)\otimes\mathcal H_W
\]
the twisted transport is defined\par
\[
\widetilde U_k(e_z\otimes v)
=
\sum_{\tau_{k+1,k}z'=z}
w_k(z'|z)\,
e_{z'}\otimes
\mathcal C_{\mathrm{Weil}}^{\,m(z\to z')}v.
\]
In matrix form,\par
\[
[\widetilde U_k]_{z',z}
=
\mathbf 1_{\{\tau z'=z\}}
\sqrt{\frac{\mu_{k+1}(z')}{\mu_k(z)}}
\,
\mathcal C_{\mathrm{Weil}}^{\,m(z\to z')}.
\]
This formula materially preserves all descendants. It does not add distinct paths as though they were the same: ledger, orientation, carry, boundary, memory, route, and terminal remain present in the label \(z'\).\par
The isometry is proved directly. Descendants of different parents are orthogonal; for each parent, the weights have squared sum one; and every power of \(\mathcal C_{\mathrm{Weil}}\) is unitary. Therefore,\par
\[
\widetilde U_k^*\widetilde U_k=I.
\]
The composition is also exact:\par
\[
\widetilde U_{\ell,j}\widetilde U_{j,k}
=
\widetilde U_{\ell,k},
\]
because the weights telescope and the memory cocycle is additive.\par
There is even an explicit gauge equivalence. If\par
\[
G_k(e_z\otimes v)
=
e_z\otimes
\mathcal C_{\mathrm{Weil}}^{\,q_{\mathrm{mem}}(z)}v,
\]
then\par
\[
\widetilde U_k
=
G_{k+1}(U_k\otimes I)G_k^*.
\]
Therefore, the twisted system is not an arbitrary decoration: it is the native cylindrical transport viewed in the spectral gauge determined by memory.\par
Define now the measurement vector\par
\[
\Omega_k
=
\sum_{z\in Z_k}
\sqrt{\mu_k(z)}\,e_z,
\qquad
\|\Omega_k\|=1,
\]
and the fit\par
\[
J_k:\mathcal H_W\longrightarrow\mathscr K_k,
\qquad
J_kv=\Omega_k\otimes v.
\]
In a full rotation, all paths possess a memory increment of one. Therefore,\par
\[
\boxed{
\widetilde U_{\Gamma_9,k}J_k
=
J_{k+9}\mathcal C_{\mathrm{Weil}}
}
\]
and, after \(n\) turns,\par
\[
\boxed{
\widetilde U_{\Gamma_9^n,k}J_k
=
J_{k+9n}\mathcal C_{\mathrm{Weil}}^n.
}
\]
This is the announced intertwiner, constructed within the same causal chain.\par
\Needspace{7\baselineskip}
\subsection{Orbital identity under the conjugation between the \(2\) spectral realizations}
Let \(\psi_\gamma\) be a vector of the spectral fiber associated with \(\gamma\). Then\par
\[
\mathcal C_{\mathrm{Weil}}\psi_\gamma
=
\tau_\gamma\psi_\gamma.
\]
Applying the preceding square:\par
\[
\widetilde U_{\Gamma_9^n}J\psi_\gamma
=
J\mathcal C_{\mathrm{Weil}}^n\psi_\gamma
=
\tau_\gamma^nJ\psi_\gamma.
\]
But the character of the memory corresponding to \(\tau_\gamma\) is\par
\[
\chi_{\tau_\gamma}(m)
=
\tau_\gamma^m.
\]
Since\par
\[
m(\Gamma_9^n)=n,
\]
se obtiene\par
\[
\boxed{
\widetilde U_{\Gamma_9^n}J\psi_\gamma
=
\chi_{\tau_\gamma}
\!\left(m(\Gamma_9^n)\right)
J\psi_\gamma.
}
\]
This is the exact equality sought.\par
It does not mean that zero number 37 is turn number 37. It means that each zero selects a character of the complete memory group and that this character traverses all turns:\par
\[
n\longmapsto\tau_\gamma^n.
\]
Por tanto:\par
\begin{itemize}
\item the nonadic turn provides the integer variable \(n\);
\item the memory provides the group \(\mathbb Z\);
\item the zero provides a character \(\chi_{\tau_\gamma}\in\widehat{\mathbb Z}\);
\item the spectral orbit and the ledger orbit coincide term by term after the embedding \(J\).
\end{itemize}
The preceding caution is overcome without falling into “nine zeros per cycle.” There is no zero--turn enumeration. There is a harmonic decomposition of the same turn dynamics.\par
\Needspace{7\baselineskip}
\subsection{Circular compactification and spiral lifting of the spectral orbit}
A crucial distinction must be preserved. The phase is unitary:\par
\[
\mathcal C_{\mathrm{Weil}}^*
\mathcal C_{\mathrm{Weil}}=I.
\]
The strict dynamics of the loss sheet is\par
\[
M_9=qV_9,
\qquad
q=\frac59.
\]
The correct set operator is\par
\[
\widetilde M_9
=
\frac59\,\widetilde U_{\Gamma_9}.
\]
On the spectral fiber,\par
\[
\widetilde M_9^nJ\psi_\gamma
=
\left(\frac59\right)^n
\tau_\gamma^n
J\psi_\gamma.
\]
Thus two images appear simultaneously:\par
\begin{itemize}
\item at the normalized boundary, the orbit is circular:
\end{itemize}
  \[
  \tau_\gamma^n\in S^1;
  \]
\begin{itemize}
\item in the strict realization, it is a spiral:
\end{itemize}
  \[
  \left(\frac59\right)^n\tau_\gamma^n\longrightarrow0.
  \]
The terminal zero does not erase the phase or the genealogy. The radius is extinguished, but the ledger preserves how many turns were performed and which character was transported.\par
Telescopy is exact:\par
\[
I
=
\sum_{r=0}^{L-1}
\widetilde M_9^{*r}
D_9^*D_9
\widetilde M_9^r
+
\widetilde M_9^{*L}\widetilde M_9^L,
\]
with\par
\[
D_9^*D_9
=
I-\widetilde M_9^*\widetilde M_9
=
\frac{56}{81}I,
\]
y terminal\par
\[
\widetilde M_9^{*L}\widetilde M_9^L
=
\left(\frac{25}{81}\right)^LI
\longrightarrow0.
\]
This preemptively answers a type objection: a contraction has not been conflated with a unitary. The phase resides in \(\mathcal C_{\mathrm{Weil}}\); the radial loss resides in \(5/9\).\par
\Needspace{7\baselineskip}
\subsection{Compatibility between genealogical cylinders, generation of constants, and spectral entangler}
Your intuition about cylinders is correct, with a precise distinction.\par
Numerical cylinders are half-open:\par
\[
C(N_K,6K)
=
\left[
\frac{N_K}{3^{6K}},
\frac{N_K+1}{3^{6K}}
\right),
\]
and each one is divided into 729 subcylinders. Their diameter satisfies\par
\[
\operatorname{diam}C(N_K,6K)=3^{-6K}\longrightarrow0.
\]
The genealogical cylinders \([a]_k\) of the inverse limit are clopen balls of the ultrametric:\par
\[
d_\vartheta(x,y)
=
\vartheta^{N(x,y)}.
\]
Therefore, "radius tending to zero" has two publications:\par
\begin{enumerate}
\item In the Archimedean reading, the intersection of nested intervals publishes a unique real number.
\item In the enriched state, the complete history continues to preserve all prefixes that produced that real number.
\end{enumerate}
Thus, in accordance with the \hyperref[sec:tesis-inaugural-helicoidal-hmt-md]{tesis inaugural}, the same coinductive system publishes \(\pi\), \(e\), \(\varphi\), and \(\alpha\) correlatively. In the internal operator order, the register \(U_{12}^{\mathrm{sgn}}\), the seal \(K\), and dodecaphasic carry act once \(P,E,\Phi\) are available: they publish and certify coordinate \(\alpha\) of that same correlated coinductive generation, without incorporating it as a causally external object. Archimedean sections and signed publication retain distinct codomains without breaking their shared genealogy.\par
These constants do not select the Riemann phases. Their relevance is deeper: they show that the same system of finite cylinders can produce a continuous object without destroying the history that generated it.\par
The new intertwiner uses exactly this property. The spectral circle is not superimposed on the continuum from outside: it appears as the dual of its integer memory.\par
\Needspace{7\baselineskip}
\subsection{Graded invariant of the spectral corridor \(4\to2\to6\to54\)}
The cadena \(4\to54\) can now situarse without numerologia.\par
The central block of APP is\par
\[
B_c
=
\begin{pmatrix}
7&2\\
2&7
\end{pmatrix}
=
9P_++5P_-.
\]
Its eigenvalues are \(9\) and \(5\), so that\par
\[
9-5=4,
\qquad
q=\frac59.
\]
The typed hierarchy is\par
\[
\underbrace{9-5}_{4};
\qquad
\underbrace{2}_{\text{coupling}};
\qquad
\underbrace{51-45}_{6};
\qquad
\underbrace{9\cdot6}_{54}.
\]
In the connection,\par
\[
\Gamma_{54}=\Gamma_9^6.
\]
By the new entangler,\par
\[
\widetilde U_{\Gamma_{54}}J
=
J\mathcal C_{\mathrm{Weil}}^6,
\]
and, in the fiber \(\gamma\),\par
\[
\widetilde U_{\Gamma_{54}}J\psi_\gamma
=
\tau_\gamma^6J\psi_\gamma.
\]
In the strict realization,\par
\[
\widetilde M_{54}J\psi_\gamma
=
\left(\frac59\right)^6
\tau_\gamma^6J\psi_\gamma.
\]
This is the exact spectral incidence of depth 54: six nonadic turns, six units of memory, and sixth power of the character.\par
It must be distinguished from other objects whose value is also 54:\par
\begin{itemize}
\item the residual APP defect;
\item the trace oriented APP–Fock;
\item the norm of the radial vector of \(A_2^{12}\);
\item the coefficient of the character \(t_3\);
\item the value
\end{itemize}
  \[
  54=\operatorname{Tr}(3B\mid V^\natural_2).
  \]
That all being 54 has structural importance, but it does not automatically make them the same morphism.\par
\Needspace{7\baselineskip}
\subsection{Leech, \(V^\natural\) and the Monster}
The exceptional chain remains exact:\par
\[
A_2^{12}
\longrightarrow
C_W=[12,6,6]_3
\longrightarrow
N_\alpha\cong\Lambda_{24}
\longrightarrow
V^\natural
\longrightarrow
\mathbb M.
\]
Transverse rigidity selects\par
\[
m=12,
\qquad
F(12)=R(12)=54.
\]
The free-fixed isometry of order three produces the orbifold\par
\[
\operatorname{Orb}_{C_3}
(V_{N_\alpha},\widehat g_c)
\cong V^\natural,
\]
and the dual automorphism belongs to the class \(3B\):\par
\[
T_{3B}
=
q^{-1}+54q+\cdots.
\]
Now, the Monster does not generate the phases \(\tau_\gamma\). It is finite, whereas the memory characters range over all of \(S^1\). In particular, an element \(3B\) has order three, but a general spectral phase is not a cube root of unity.\par
The correct meeting is a product of representations:\par
\[
\mathscr K^{\mathrm{gold}}
=
\mathscr H_{\mathrm{cyl}}
\otimes
\mathcal H_W
\otimes
\mathcal H_{V^\natural}.
\]
Memory and the Cayley structure act on the first two factors; the Monster acts on the third:\par
\[
V_{\mathrm{mem}}
=
\widetilde U_{\Gamma_9}\otimes I,
\qquad
\rho_{\mathbb M}(g)
=
I\otimes I\otimes\rho_{V^\natural}(g).
\]
Both actions commute. Thus, each spectral orbit can simultaneously preserve:\par
\begin{itemize}
\item su altura \(\gamma\);
\item its phase \(\tau_\gamma\);
\item its nonadic memory \(n\);
\item its spectral multiplicity;
\item its transverse Monster type.
\end{itemize}
But the multiplicity of a zero is not identified by decree with a monstrous dimension. That identification would require an additional intertwiner between the spectral fiber and a concrete representation of the Monster.\par
If a genuine computable family of entangled traces remains open:\par
\[
\Theta_g(n;a,\beta)
=
\operatorname{Tr}
\left(
e^{-aH_W^2}\mathcal C_{\mathrm{Weil}}^n
\right)
\operatorname{Tr}_{V^\natural}
\left(
g\,e^{-\beta(L_0-1)}
\right).
\]
After spectral recognition,\par
\[
\Theta_g(n;a,\beta)
=
\left(
\sum_\gamma
m_\gamma e^{-a\gamma^2}\tau_\gamma^n
\right)
T_g(e^{-\beta}).
\]
This indeed is a new rigorously typed RH–Moonshine path: moments of the orbits of the zeros, nonadic memory, and monstrous characters in a single regularized observable.\par
\Needspace{7\baselineskip}
\subsection{Finite Approximation of the Spectral Field via Matrix Truncations}
The construction possesses a finite natural matrix.\par
For any unitary \(C\), define the nine-phase stitching\par
\[
\mathbb U_{C,9}
=
\sum_{r=1}^{8}
|r+1\rangle\langle r|\otimes I
+
|1\rangle\langle9|\otimes C.
\]
Then\par
\[
\mathbb U_{C,9}^9
=
I\otimes C.
\]
Con\par
\[
u_9=\frac1{3}(1,\ldots,1),
\qquad
\iota_9v=u_9\otimes v,
\]
visible compression is\par
\[
A_{C,9}
=
\iota_9^*\mathbb U_{C,9}\iota_9
=
\frac{8I+C}{9},
\]
and the defect satisfies\par
\[
c_{C,9}^*c_{C,9}
=
\frac{8}{81}
(I-C)^*(I-C).
\]
Taking \(C=\mathcal C_{\mathrm{Weil}}\) yields a nine-phase matrix constructed without introducing zeros.\par
To make the Weil fiber finite as well, test functions \(f_1,\ldots,f_N\) are
chosen and the matrices are computed from the prime--gamma--polar channels\par
\[
G_N
=
\bigl(
\mathscr I_{\mathrm{GW}}(f_i,f_j)
\bigr)_{i,j},
\]
y\par
\[
A_N
=
\bigl(
\langle H_Wf_j,f_i\rangle_W
\bigr)_{i,j}.
\]
After taking the quotient by \(\ker G_N\), the generalized problem\par
\[
A_Nv=\gamma_NG_Nv
\]
produce the spectral approximation, and its finite Cayley is\par
\[
C_N
=
\left(
G_N^{-1}A_N+\frac{i}{2}I
\right)
\left(
G_N^{-1}A_N-\frac{i}{2}I
\right)^{-1}.
\]
Its eigenvalues \(\tau_N\in S^1\) recover\par
\[
\gamma_N
=
\frac12
\cot\left(
\frac{\arg\tau_N}{2}
\right).
\]
The total cylinder matrix is then\par
\[
[\widetilde U_{k,N}]_{z',z}
=
\mathbf1_{\{\tau z'=z\}}
\sqrt{\frac{\mu_{k+1}(z')}{\mu_k(z)}}
\,C_N^{\,m(z\to z')}.
\]
This converts the conceptual approach into a specific numerical program:\par
\begin{enumerate}
\item construct \(G_N\) and \(A_N\) without introducing zeros;
\item form \(C_N\);
\item insert it into the nonadic seam;
\item diagonalizing the phases;
\item recover the heights;
\item verify stability under simultaneous increase of \(N\) and cylindrical depth.
\end{enumerate}
Rigorous convergence requires proving that the finite subspaces form a core for \(H_W\) and that the approximations converge in the strong resolvent sense. This is a subsequent numerical certificate; it does not affect the infinite intertwiner, which has already been proved.\par
\Needspace{7\baselineskip}
\subsection{Independence of the spectral intertwiner with respect to zeros, contractions, route multiplicities, and transverse symmetries}
\textbf{“The zeros were introduced to fabricate the map.”} No. \(\mathcal H_W\), \(H_W\), their resolvent, and \(\mathcal C_{\mathrm{Weil}}\) are constructed from the prime–gamma–polar form. The zeros appear afterward to recognize the spectrum.\par
\textbf{“A contraction has been identified with a unitary operator.”} No. The phase is \(\mathcal C_{\mathrm{Weil}}\); the contraction is \((5/9)\mathcal C_{\mathrm{Weil}}\).\par
\textbf{“Distinct paths have been erased.”} No. Each descendant preserves its complete label in \(e_{z'}\); the orthogonality of the ledger prevents two routes from being confused.\par
\textbf{“It is claimed that there are nine zeros per cycle.”} No. Each zero determines a character that is evaluated over all turns.\par
\textbf{“The Monster produces the heights.”} No. The phase circle arises from \(\widehat{\mathbb Z_{\mathrm{mem}}}\). The Monster organizes transverse symmetry and commutes with that action.\par
\textbf{“The different 54s are the same object.”} No. Defect, depth, norm, and character value are kept separate. Their coincidences remain connected by maps, not by bare numerical equality.\par
\textbf{“Only a pre-existing operator \(V_9\) has been renamed.”} No. A specific explicit realization of the codomain of \(\Gamma_9\) has been constructed. To identify it with any pre-existing ambient representative of \(V_9\), specified only up to unitary equivalence, the preceding gauge must be chosen compatibly. On the relevant cyclic subspace, the unitary class is already determined.\par
\Needspace{7\baselineskip}
\subsection{Domain of the Cayley--Pontryagin conjugation and theorem of the double circle}
APP–TRIT–TPK, the cylinders, the projective measure, the memory cocycle, \(\Gamma_9\), the contraction \(5/9\), the chain \(4\to54\), the constants, and the exceptional corridor constitute the initial structure. On that domain, the cylindrical transport twisted by \(\mathcal C_{\mathrm{Weil}}^{m(\gamma)}\) is constructed. The double-circle theorem establishes the exact intertwining
  \[
  \widetilde U_{\Gamma_9}J
  =
  J\mathcal C_{\mathrm{Weil}}.
  \]
Certification of convergence of the finite matrices and the unfactorized study of the RH–Moonshine traces constitute the open boundary beyond the theorem.
Therefore, the strong conclusion is:\par
\[
\boxed{
\text{zeros are not the turns;
they are the spectral characters of the memory of the turns.}
}
\]
And, after the canonical matching,\par
\[
\boxed{
\text{the spectral orbit of the zero and the ledger orbit are literally the same orbit.}
}
\]

\section{Effective reconstruction of the Weil spectral field}
\Needspace{5\baselineskip}
\subsection{Construction of the spectral field independent of previous data on the zeros}
The compactly supported family is adopted.\par
\[
\beta_A(x)=Z_A^{-1}
\exp\!\left[-\frac{1}{1-(x/A)^2}\right]
\mathbf 1_{\{|x|<A\}},
\qquad
f_{\omega,A}(x)=\beta_A(x)e^{-i\omega x}.
\]
The modulation \(\omega\) was not chosen using zeros. The interval was swept in advance\par
\[
5\leq\omega\leq50,
\qquad
\Delta\omega=0.025.
\]
For each \(\omega\) the calculation is performed\par
\[
M_A(\omega)
=
\mathscr I_{\rm GW}
(f_{\omega,A},f_{\omega,A})
\]
via the formula constructed from the arithmetic data\par
\[
\begin{aligned}
\mathscr I_{\rm GW}(f,g)
={}&c_\Gamma C_{f,g}(0)
+\int_0^\infty
\mu_\Gamma(a)
\bigl[
2C_{f,g}(0)-C_{f,g}(a)-C_{f,g}(-a)
\bigr]\,da\\
&-\sum_{p,k}
\frac{\log p}{p^{k/2}}
\bigl[
C_{f,g}(k\log p)+C_{f,g}(-k\log p)
\bigr]\\
&+P_f^+\overline{P_g^-}
+P_f^-\overline{P_g^+},
\end{aligned}
\]
where\par
\[
C_{f,g}(a)=\int_{\mathbb R}f(x)\overline{g(x-a)}\,dx,
\]
\[
\mu_\Gamma(a)=\frac{e^{-a/2}}{1-e^{-2a}},
\qquad
c_\Gamma=-\gamma_{\!E}-\frac{\pi}{2}-3\log2-\log\pi,
\]
y\par
\[
P_f^\pm=\int_{\mathbb R}f(x)e^{\pm x/2}\,dx.
\]
Neither \(\zeta(s)\), a zero-finding routine, Riemann–Siegel, Odlyzko, nor a table of heights was involved.\par
With \(A=3\), the first blind sweep produced ten dominant maxima approximately at\par
\[
14.125,\ 21.025,\ 25.000,\ 30.425,\ 32.925,\ 37.575,\ 40.925,\ 43.325,\ 48.000,\ 49.775.
\]
The lateral lobes had a mass about a hundred times smaller.\par
\Needspace{5\baselineskip}
\subsection{Refinement of the independent spectral field of the target values}
I increased the compact width \(A\). This simultaneously increases the spectral resolution and the number of prime powers used:\par
\Needspace{8\baselineskip}
\begin{center}
\small
\begin{tabularx}{\textwidth}{@{}RRRRR@{}}
\toprule
\(n\) & \(A=3\) & \(A=4\) & \(A=5\) & \(A=6\) \\
\midrule
1 & 14.134699719 & 14.134721825 & 14.134724556 & 14.134725023 \\
2 & 21.022713497 & 21.022092670 & 21.022033480 & 21.022034382 \\
3 & 25.010192460 & 25.010816662 & 25.010862963 & 25.010863775 \\
4 & 30.431966224 & 30.424255364 & 30.424903280 & 30.424992486 \\
5 & 32.927889053 & 32.935726409 & 32.935037177 & 32.934942801 \\
6 & 37.585309844 & 37.585911673 & 37.586134491 & 37.586172107 \\
7 & 40.926249022 & 40.917939546 & 40.919104383 & 40.918785046 \\
8 & 43.320460106 & 43.328124621 & 43.326729729 & 43.327013819 \\
9 & 47.996518745 & 48.009530251 & 48.005271807 & 48.005312108 \\
10 & 49.781332020 & 49.769281661 & 49.773747976 & 49.773697939 \\
\bottomrule
\end{tabularx}
\end{center}
In \(A=6\) exactly 15.040 prime powers \(p^k\) entered, with \(p^k\leq e^{12}\). For the first maximum, the numerical ledger was:\par
\[
\begin{array}{lr}
\text{prime channel} & 8.077017879363272,\\
\text{gamma integral} & 6.182517531024263,\\
c_\Gamma I & -5.372183419225665,\\
\text{polar channel} & -1.0505899\times10^{-8},\\ \hline
\text{total mass} & 8.887351980655973.
\end{array}
\]
The gamma tail beyond \(a=65\) was bounded by\par
\[
1.54\times10^{-14}.
\]
This reveals a conceptual point: zero is not inserted into the calculation.
It appears as the stable center at which the four components of the
functional close.\par
\Needspace{5\baselineskip}
\subsection{Finite matrix approximants of the spectral operator}
An isolated Gram matrix does not contain heights in identifiable form. Its eigenvalues depend on the basis and represent masses or energies, not spectral positions.\par
Therefore, three matrices are constructed:\par
\[
G_{ij}=\mathscr I_{\rm GW}(f_j,f_i),
\]
\[
K_{ij}=\mathscr I_{\rm GW}(H_Wf_j,f_i),
\]
\[
Q_{ij}=\mathscr I_{\rm GW}(H_Wf_j,H_Wf_i).
\]
With the positive orientation used in calculus,\par
\[
H_W=i\,\partial_x
\]
on the reflected family \(f_{\omega,A}=\beta_Ae^{-i\omega x}\). The conjugate orientation \(-i\partial_x\) returns the opposite spectral sheet; both produce the same positive heights by the symmetry \(\gamma\leftrightarrow-\gamma\).\par
This generator is not introduced arbitrarily. The established positivity permits construction of the GNS representation of \(\mathscr I_{\rm GW}\). The translations are strongly continuous isometries of that GNS representation, and Stone's theorem produces their self-adjoint generator. On the test core, the derivation of the translations is precisely \(i\partial_x\).\par
The three matrices were computed by means of the same prime–gamma–polar evaluator. Neither \(K\) nor \(Q\) was constructed from a spectral list.\par
In block \(10\times10\), the hermiticity defects before numerical symmetrization were\par
\[
\|G-G^*\|_2=2.11\times10^{-14},
\]
\[
\|K-K^*\|_2=7.14\times10^{-13},
\]
\[
\|Q-Q^*\|_2=1.84\times10^{-10}.
\]
The eigenvalues of \(G\) were\par
\[
\begin{aligned}
4.90806959,\ 5.50537820,\ 5.61057540,\ 5.88666432,\ 5.91920889,\\
5.93049827,\ 5.96524271,\ 6.27603538,\ 6.33591428,\ 7.02999383.
\end{aligned}
\]
All are positive and the block condition is only\par
\[
\kappa(G)=1.43233.
\]
There was no need to hide a negative eigenvalue or artificially project an uncomfortable part.\par
The metric was whitened:\par
\[
\widehat H_N
=
G_N^{-1/2}K_NG_N^{-1/2},
\]
and the hermitic pencil was resolved\par
\[
K_Nv=\gamma_NG_Nv.
\]
This is the point that converts gram masses into heights.\par
\Needspace{5\baselineskip}
\subsection{Cayley Transform and Initial Circular Compactification of the Spectrum}
To each \(\widehat H_N\) was applied\par
\[
C_N=
\left(\widehat H_N+\frac{i}{2}I\right)
\left(\widehat H_N-\frac{i}{2}I\right)^{-1}.
\]
If \(\gamma\) is an eigenvalue of the generator, its phase is\par
\[
\tau_\gamma
=
\frac{\gamma+i/2}{\gamma-i/2},
\qquad |\tau_\gamma|=1.
\]
The inverse recovery is\par
\[
\gamma
=
\frac{i}{2}\frac{\tau+1}{\tau-1}
=
\frac12\cot\!\left(\frac{\arg\tau}{2}\right).
\]
The calculated unitarity defect was\par
\[
\|C_N^*C_N-I\|_2=1.45\times10^{-15}.
\]
The first three phases were\par
\[
\tau_1=
0.997500505583064
+0.070659333152323\,i,
\]
\[
\tau_2=
0.998869228927548
+0.047542228615055\,i,
\]
\[
\tau_3=
0.999201013749813
+0.039966662624574\,i.
\]
This is the first circle: the spectral line is compactified without losing any height, because the Cayley map is bijective between \(\mathbb R\) and the circle with the terminal point removed.\par
\Needspace{5\baselineskip}
\subsection{Exact immersion of the spectrum into the nonadic circle}
Once \(C_N\) has been constructed, the nine-position sewing matrix is defined\par
\[
S_{N,9}
=
\sum_{r=1}^{8}
|r+1\rangle\langle r|\otimes I
+
|1\rangle\langle9|\otimes C_N.
\]
A zero is not assigned to each gate. The nine positions are the transport calendar. The spectral phase is applied only when the traversal completes the turn.\par
The exact identity is\par
\[
S_{N,9}^{\,9}=I_9\otimes C_N.
\]
Therefore, if \(J_Nv=u_9\otimes v\), the full-turn holonomy satisfies\par
\[
S_{N,9}^{\,9}J_N=J_NC_N.
\]
This is the entanglement of the double circle:\par
\begin{itemize}
\item the Cayley circle contains the spectral phase;
\item the nonadic circle transports the state for nine steps;
\item a complete turn applies the Cayley phase exactly once;
\item memory records the number of revolutions;
\item no zero is identified with an isolated gate.
\end{itemize}
For the ten candidates a matrix \(90\times90\) was constructed. The controls were:\par
\[
\|S_{N,9}^*S_{N,9}-I\|_2
=
4.44\times10^{-16},
\]
\[
\|S_{N,9}^{\,9}-I_9\otimes C_N\|_2=0.
\]
The visible contraction is\par
\[
M_N=\frac59S_{N,9},
\qquad
D_N=\frac{\sqrt{56}}9I.
\]
It was verified\par
\[
M_N^*M_N=\frac{25}{81}I,
\qquad
D_N^*D_N=\frac{56}{81}I,
\]
and the telescopy\par
\[
I=
\sum_{r=0}^{L-1}
M_N^{*r}D_N^*D_NM_N^r
+
M_N^{*L}M_N^L.
\]
Para \(L=20\),\par
\[
\left\|
\sum_{r=0}^{19}M_N^{*r}D_N^*D_NM_N^r
+
M_N^{*20}M_N^{20}
-I
\right\|_2
=
3.33\times10^{-16},
\]
while the terminal was\par
\[
\|M_N^{*20}M_N^{20}\|_2
=
6.1531824951\times10^{-11},
\]
exactly equal, within machine precision, to\par
\[
\left(\frac{25}{81}\right)^{20}.
\]
The visible loss does not destroy the phase: the phase resides in the unitary dilation, and the factor \(5/9\) regulates how much remains visible on each turn.\par
\Needspace{5\baselineskip}
\subsection{Recovery of Heights and Residues with Posterior External Contrast}
For each generalized vector \(v\), the residue was computed without zeros through\par
\[
r_N(\gamma,v)^2
=
\frac{
v^*
\left(
Q_N-2\gamma K_N+\gamma^2G_N
\right)v
}{
v^*G_Nv
}.
\]
This residual measures how far the finite vector lies from a genuine spectral vector. A height is not accepted merely because the pencil produces a number.\par
\Needspace{8\baselineskip}
\begin{center}
\small
\begin{tabularx}{\textwidth}{@{}RRRR@{}}
\toprule
\(n\) & altura producida & residue & error in the posterior contrast \\
\midrule
1 & 14.134725141525 & \(8.15\times10^{-5}\) & \(2.10\times10^{-10}\) \\
2 & 21.022039638825 & \(4.82\times10^{-5}\) & \(5.34\times10^{-11}\) \\
3 & 25.010857580595 & \(1.15\times10^{-4}\) & \(4.49\times10^{-10}\) \\
4 & 30.424876128079 & \(2.32\times10^{-4}\) & \(2.22\times10^{-9}\) \\
5 & 32.935061595211 & \(3.98\times10^{-4}\) & \(7.47\times10^{-9}\) \\
6 & 37.586178242423 & \(1.18\times10^{-3}\) & \(8.36\times10^{-8}\) \\
7 & 40.918719778888 & \(3.14\times10^{-3}\) & \(7.67\times10^{-7}\) \\
8 & 43.327078423026 & \(7.18\times10^{-3}\) & \(5.14\times10^{-6}\) \\
9 & 48.005155742413 & \(6.22\times10^{-3}\) & \(4.86\times10^{-6}\) \\
10 & 49.773888321513 & \(1.45\times10^{-2}\) & \(5.58\times10^{-5}\) \\
\bottomrule
\end{tabularx}
\end{center}
The final column was computed only after the outputs had been frozen. It played no role in the sweep, in \(G\), in \(K\), in \(Q\), in the Cayley transform, or in the nonadic seam.\par
\Needspace{5\baselineskip}
\subsection{Convergence and preservation of the spectral genealogical record}
After the Riemann–Weil resolution, the functional possesses the representation\par
\[
\mathscr I_{\rm GW}(f,g)
=
\sum_\gamma
m_\gamma
\widehat f(-\gamma)
\overline{\widehat g(-\gamma)}.
\]
This is not used to calculate the entries; it serves to recognize which spectrum we have just obtained.\par
The GNS remains identificado with the espacio \(L^2\) of the measure atomica\par
\[
\nu_W=\sum_\gamma m_\gamma\delta_\gamma.
\]
The generator of translations acts by multiplication by \(\gamma\). Therefore:\par
\begin{itemize}
\item \(G\) contains the scalar products;
\item \(K\) inserts a power of \(\gamma\);
\item \(Q\) inserta \(\gamma^2\);
\item the pencil \(K-\gamma G\) localizes the atoms;
\item Cayley transforms them into unitary phases;
\item the nonadic seam transports those phases through the ledger.
\end{itemize}
The resolution is not being used to prove RH again with ten matrices. The already obtained resolution is used to convert the explicit formula into a spectral-extraction algorithm that does not receive the zeros as input.\par
\Needspace{5\baselineskip}
\subsection{Convergence of the reconstructed spectral field and certifiable extension}
The computational cycle is now closed:\par
\[
\text{primes+gamma+polar}
\longrightarrow
(G,K,Q)
\longrightarrow
\widehat H_N
\longrightarrow
C_N
\longrightarrow
S_{N,9}
\longrightarrow
\{\gamma_n\}.
\]
The literal interpretation of the closure circle is also constructed:\par
\[
\text{one complete nonadic turn}
\quad\Longleftrightarrow\quad
\text{one application of the spectral phase},
\]
but not\par
\[
\text{a nonadic phase}
\quad\Longleftrightarrow\quad
\text{a zero}.
\]
Each zero defines a character of the ledger,\par
\[
n\longmapsto\tau_\gamma^n,
\]
and in the contractive publication it appears as\par
\[
n\longmapsto
\left(\frac59\right)^n\tau_\gamma^n.
\]
This is the concrete orbit: not a repetition of heights, but the evolution of a single spectral phase through the memory of turns.\par
Nor is the number \(54\) from the exceptional branch conflated with a number of zeros or turns. The \(54\) of APP–Witt–Leech–Monster is a defect, a norm, and a graded character in its domain. It may coexist with the spectral circle through a product of representations, but it neither selected nor normalized any of the heights calculated here.\par
Mathematical convergence rests on a cofinal family of windows and compact modes constituting a graph core for the GNS generator. The calculation performed supplies four window refinements and a posteriori residuals. A certified version using interval arithmetic could convert every decimal into a guaranteed interval, but no conceptual or operatorial map is any longer missing: this would be the numerical strengthening of a chain already constructed from beginning to end.\par
The explicit formula and the joint nonadic connection constitute the initial mathematical structure. The matrices \(G,K,Q\), the Cayley–Galerkin operator, and the finite double-circle intertwining constitute the numerical construction. The ten heights are extracted directly from the arithmetic form and constitute the result of that calculation.\par
